Dans un lycée, des élèves de première souhaitent installer des vélos générateurs d'électricité dans la cafétéria pour recharger des téléphones portables.

	À l'aide des quatre documents ci-dessous, indiquer si chacune des affirmations suivantes est vraie ou fausse. \textbf{Chaque réponse devra être justifiée}.

	\bigskip

	\textbf{Affirmation 1 :} Le budget du lycée permet de financer cinq vélos générateurs d'électricité (vélo, chargeur et livraison).

	\bigskip

	\textbf{Affirmation 2 :} Il faut pédaler pendant 24 minutes à intensité modérée pour recharger un téléphone portable à l'aide d'un vélo générateur d'électricité.

	\bigskip

	\textbf{Affirmation 3 :} On choisit au hasard un élève du lycée.

	La probabilité que l'élève choisi soit un élève de première en sachant qu'il est favorable au projet est supérieure à un quart.

	\bigskip

	\textbf{Affirmation 4 :} Le projet des élèves de première va pouvoir être lancé.

	\subsection*{Documents}
	\textbf{Document 1 : Fiche d'un vélo générateur d'électricité}

	\begin{center}
	\fbox{\begin{minipage}{0.9\linewidth}
		\textbf{Tarifs pour l'installation de vélos générateurs d'électricité}
		\begin{itemize}[left=5mm]
			\item Coût d'un vélo sans chargeur : 300 €
			\item Coût du chargeur, pour chaque vélo : 45 €
			\item Coût de livraison par vélo : 60 €
		\end{itemize}

		\textbf{Énergie produite par un vélo générateur d'électricité}

		Lorsqu'une personne pédale à intensité modérée pendant une heure, le dispositif permet de récupérer 100 kJ (kilojoules).
	\end{minipage}}
	\end{center}

	\textbf{Document 2 : Énergie nécessaire pour charger quelques appareils}
	\begin{center}
		\begin{tblr}{width=0.9\linewidth,colspec={X[c,2.5]*4{X[c]}},hlines, vlines, column{1} = {bg=gray!30},rows={m}}
			Appareils électroniques & Ordinateur portable & Casque audio & Téléphone portable & Montre connectée \\
			Énergie moyenne nécessaire pour une charge complète
			 & 200 kJ & 18 kJ & 40 kJ & 5 kJ \\
		\end{tblr}
	\end{center}

	\needspace{4\baselineskip}
	\textbf{Document 3 : Résultats du sondage réalisé auprès des élèves du lycée sur les vélos générateurs d'électricité}
	\begin{center}
		\begin{tblr}{width=0.9\linewidth,colspec={*4{X[c,m]}},column{1}={bg=gray!30},row{1}={bg=gray!30},cell{1}{1}={bg=white},hline{1}={2-Z}{solid},hline{2-Z}={solid},vline{1}={2-Z}{solid},vline{2-Z}={solid}}
			          & Favorable & Défavorable & Total \\
			Seconde   & 260       & 170         & 430   \\
			Première  & 180       & 120         & 300   \\
			Terminale & 120       & 150         & 270   \\
			Total     & 560       & 440         & 1000
		\end{tblr}
	\end{center}

	\textbf{Document 4 : Message de la proviseure du lycée aux élèves porteurs du projet}

	\begin{center}
	\fbox{\begin{minipage}{0.9\linewidth}
	\textbf{Objet : Lancement du projet "vélos générateurs d'électricité" -- conditions à
	remplir}

	Chers élèves de première,

	Ce projet ne pourra être lancé que si les \textbf{deux critères suivants sont remplis} :

	\begin{itemize}
		\item \textbf{plus de la moitié des élèves du lycée} doivent se déclarer favorables à cette initiative ;

		\item \textbf{au moins $60\,\%$ des élèves de première} doivent également y être \textbf{favorables}.

		Dans le cas où le projet serait validé, \textbf{le budget du lycée sera fixé à \np[\euro]{2000}} pour financer l'achat des vélos et des chargeurs ainsi que la livraison des vélos.
	\end{itemize}

	Bien cordialement,

	La Proviseure
	\end{minipage}}
	\end{center}
